\chapter{Calcul intégral, encadrements et suites d’intégrales}
\section*{Exercice 20.1 -- Intégration par parties}
[Méthode] Calculer $\int_0^1xe^x\,dx$.

\section*{Exercice 20.2 -- Aire entre deux courbes}
[Méthode] Calculer l'aire comprise entre $y=x$ et $y=x^2$ de 0 à 1.

\section*{Exercice 20.3 -- Suite d’intégrales}
[Preuve] Pour $I_n=\int_0^1x^ne^{-x}\,dx$, prouver que $I_n\to0$ et établir $I_n=nI_{n-1}-e^{-1}$ pour $n\ge1$.

\section*{Exercice 20.4 -- Somme de Riemann}
[Approfondissement] Calculer $\lim_{n\to\infty}\frac1n\sum_{k=1}^n\frac1{1+k/n}$.

\section*{Exercice 20.5 -- Trapèzes et convexité}
[Raisonnement] Pour une fonction convexe sur $[a,b]$, montrer $\int_a^bf\le(b-a)(f(a)+f(b))/2$.

\section*{Exercice 20.6 -- Série harmonique}
[Preuve] Utiliser une intégrale pour montrer que $H_n=\sum_{k=1}^n1/k$ diverge.
